\documentclass[10pt,a4paper]{amsart}
\usepackage{amsmath,amssymb,amsthm,graphicx}
% Retained preamble for the Stacks Project chapter snapshot in this directory.
% The chapter's own preamble.tex is a copy of the shared Stacks Project preamble; the block
% below is sliced MECHANICALLY, line for line, from it. Nothing was added, removed or
% redefined; the wrapper class, the xr-hyper/hyperref wiring and the externaldocument list
% are omitted because the assembler supplies the document class.
%
% lines taken from preamble.tex: 3, 5, 9, 11, 12, 156, 157, 158, ... (51 lines total)
%
% package coverage restored from the chapter's own preamble (2026-10-13): xy, xy-2cell

\usepackage[all]{xy}
\xyoption{2cell}
\UseAllTwocells

\newenvironment{reference}{\comment}{\endcomment}
\newenvironment{slogan}{\comment}{\endcomment}
\newenvironment{history}{\comment}{\endcomment}
\theoremstyle{plain}
\newtheorem{theorem}[subsection]{Theorem}
\newtheorem{proposition}[subsection]{Proposition}
\newtheorem{lemma}[subsection]{Lemma}
\theoremstyle{definition}
\newtheorem{definition}[subsection]{Definition}
\newtheorem{example}[subsection]{Example}
\newtheorem{exercise}[subsection]{Exercise}
\newtheorem{situation}[subsection]{Situation}
\theoremstyle{remark}
\newtheorem{remark}[subsection]{Remark}
\newtheorem{remarks}[subsection]{Remarks}
\numberwithin{equation}{subsection}
\def\lim{\mathop{\mathrm{lim}}\nolimits}
\def\colim{\mathop{\mathrm{colim}}\nolimits}
\def\Spec{\mathop{\mathrm{Spec}}}
\def\Hom{\mathop{\mathrm{Hom}}\nolimits}
\def\Ext{\mathop{\mathrm{Ext}}\nolimits}
\def\SheafHom{\mathop{\mathcal{H}\!\mathit{om}}\nolimits}
\def\SheafExt{\mathop{\mathcal{E}\!\mathit{xt}}\nolimits}
\def\Sch{\mathit{Sch}}
\def\Mor{\mathop{\mathrm{Mor}}\nolimits}
\def\Ob{\mathop{\mathrm{Ob}}\nolimits}
\def\Sh{\mathop{\mathit{Sh}}\nolimits}
\def\NL{\mathop{N\!L}\nolimits}
\def\CH{\mathop{\mathrm{CH}}\nolimits}
\def\proetale{{pro\text{-}\acute{e}tale}}
\def\etale{{\acute{e}tale}}
\def\QCoh{\mathit{QCoh}}
\def\Ker{\mathop{\mathrm{Ker}}}
\def\Im{\mathop{\mathrm{Im}}}
\def\Coker{\mathop{\mathrm{Coker}}}
\def\Coim{\mathop{\mathrm{Coim}}}
\DeclareMathSymbol{\boxtimes}{\mathbin}{AMSa}{"02}
\def\QCohstack{\mathcal{QC}\!\mathit{oh}}
\def\Cohstack{\mathcal{C}\!\mathit{oh}}
\def\Spacesstack{\mathcal{S}\!\mathit{paces}}
\def\Quotfunctor{\mathrm{Quot}}
\def\Hilbfunctor{\mathrm{Hilb}}
\def\Curvesstack{\mathcal{C}\!\mathit{urves}}
\def\Polarizedstack{\mathcal{P}\!\mathit{olarized}}
\def\Complexesstack{\mathcal{C}\!\mathit{omplexes}}
\def\Pic{\mathop{\mathrm{Pic}}\nolimits}
\def\Picardstack{\mathcal{P}\!\mathit{ic}}
\def\Picardfunctor{\mathrm{Pic}}
\def\Deformationcategory{\mathcal{D}\!\mathit{ef}}
\usepackage{cleveref}
\crefname{originallemma}{original lemma}{original lemmas}
\Crefname{originallemma}{Original lemma}{Original lemmas}
\makeatletter
\newlabel{algebra-lemma-characterize-smooth-over-field}{{00TT}{source}{}{sourcepointer-1}{}}
\newlabel{algebra-lemma-characterize-smooth-over-field@cref}{{[originallemma][0][]00TT}{[1][0][]source}{}{}{}}
\makeatother
\crefname{originallemma}{original lemma}{original lemmas}
\Crefname{originallemma}{Original lemma}{Original lemmas}
\makeatletter
\newlabel{algebra-lemma-differentials-finitely-presented}{{00RY}{source}{}{sourcepointer-2}{}}
\newlabel{algebra-lemma-differentials-finitely-presented@cref}{{[originallemma][0][]00RY}{[1][0][]source}{}{}{}}
\makeatother
\crefname{originallemma}{original lemma}{original lemmas}
\Crefname{originallemma}{Original lemma}{Original lemmas}
\makeatletter
\newlabel{algebra-lemma-differentials-localize}{{00RT}{source}{}{sourcepointer-3}{}}
\newlabel{algebra-lemma-differentials-localize@cref}{{[originallemma][0][]00RT}{[1][0][]source}{}{}{}}
\makeatother
\crefname{originallemma}{original lemma}{original lemmas}
\Crefname{originallemma}{Original lemma}{Original lemmas}
\makeatletter
\newlabel{algebra-lemma-finite-type-algebra-finite-nr-primes}{{0ALW}{source}{}{sourcepointer-4}{}}
\newlabel{algebra-lemma-finite-type-algebra-finite-nr-primes@cref}{{[originallemma][0][]0ALW}{[1][0][]source}{}{}{}}
\makeatother
\crefname{originallemma}{original lemma}{original lemmas}
\Crefname{originallemma}{Original lemma}{Original lemmas}
\makeatletter
\newlabel{algebra-lemma-hom-from-finitely-presented}{{0583}{source}{}{sourcepointer-5}{}}
\newlabel{algebra-lemma-hom-from-finitely-presented@cref}{{[originallemma][0][]0583}{[1][0][]source}{}{}{}}
\makeatother
\crefname{originallemma}{original lemma}{original lemmas}
\Crefname{originallemma}{Original lemma}{Original lemmas}
\makeatletter
\newlabel{algebra-lemma-localize-NL}{{00S7}{source}{}{sourcepointer-6}{}}
\newlabel{algebra-lemma-localize-NL@cref}{{[originallemma][0][]00S7}{[1][0][]source}{}{}{}}
\makeatother
\crefname{originallemma}{original lemma}{original lemmas}
\Crefname{originallemma}{Original lemma}{Original lemmas}
\makeatletter
\newlabel{algebra-lemma-simple-pieces}{{00J3}{source}{}{sourcepointer-7}{}}
\newlabel{algebra-lemma-simple-pieces@cref}{{[originallemma][0][]00J3}{[1][0][]source}{}{}{}}
\makeatother
\crefname{originallemma}{original lemma}{original lemmas}
\Crefname{originallemma}{Original lemma}{Original lemmas}
\makeatletter
\newlabel{algebra-lemma-smooth-at-point}{{07BU}{source}{}{sourcepointer-8}{}}
\newlabel{algebra-lemma-smooth-at-point@cref}{{[originallemma][0][]07BU}{[1][0][]source}{}{}{}}
\makeatother
\crefname{originallemma}{original lemma}{original lemmas}
\Crefname{originallemma}{Original lemma}{Original lemmas}
\makeatletter
\newlabel{algebra-proposition-minimal-primes-associated-primes}{{02CE}{source}{}{sourcepointer-9}{}}
\newlabel{algebra-proposition-minimal-primes-associated-primes@cref}{{[originallemma][0][]02CE}{[1][0][]source}{}{}{}}
\makeatother
\crefname{originallemma}{original lemma}{original lemmas}
\Crefname{originallemma}{Original lemma}{Original lemmas}
\makeatletter
\newlabel{algebra-theorem-nullstellensatz}{{00FV}{source}{}{sourcepointer-10}{}}
\newlabel{algebra-theorem-nullstellensatz@cref}{{[originallemma][0][]00FV}{[1][0][]source}{}{}{}}
\makeatother
\crefname{originallemma}{original lemma}{original lemmas}
\Crefname{originallemma}{Original lemma}{Original lemmas}
\makeatletter
\newlabel{lemma-rings-TI}{{0DY3}{source}{}{sourcepointer-11}{}}
\newlabel{lemma-rings-TI@cref}{{[originallemma][0][]0DY3}{[1][0][]source}{}{}{}}
\makeatother
\crefname{originallemma}{original lemma}{original lemmas}
\Crefname{originallemma}{Original lemma}{Original lemmas}
\makeatletter
\newlabel{more-algebra-lemma-pseudo-coherence-and-base-change-ext}{{087R}{source}{}{sourcepointer-12}{}}
\newlabel{more-algebra-lemma-pseudo-coherence-and-base-change-ext@cref}{{[originallemma][0][]087R}{[1][0][]source}{}{}{}}
\makeatother
\title{Stacks, Examples of Deformations: lemma-finite-type-rings-TI}
\author{The Stacks Project Authors}
\date{stacks-project:examples-defos.tex:local-snapshot-20261012; GNU FDL 1.2 or later}
\begin{document}
\maketitle
\noindent\textit{Editorial note.} Statement span is lines 1123-1132 of sources/stacks-examples-defos/pinned-original-source.tex; proof span is lines 1134-1195, the complete author proof (60 body lines).\par\medskip
\noindent\textit{Editorial note.} Navigation stubs: 12 cited-but-not-included labels are supplied as non-load-bearing pointers with their pinned tags where the snapshot tag index has them.\par\medskip
\noindent\textit{Editorial note.} extra\_packages is empty on purpose: the retained preamble already loads xy with options and a second declaration clashes in the assembler wrapper.\par\medskip
\noindent\textit{Editorial note.} Text is retained verbatim, including the author's own typos and placeholders.\par\medskip
\noindent\textit{Editorial note.} Queue index 3 of queue\_examples-defos.json (pb=2941, stmt 1123-1132, proof 1134-1195).\par\medskip
\noindent\textit{Allowed standard prerequisites (editorial).}\begin{itemize}
\item Explicit prerequisite boundary. The cited inputs are prerequisites; the argument in the retained span is not replaced by any cited result.
\end{itemize}\medskip
\ifdefined\hypertarget\hypertarget{sourcepointer-1}{}\fi
\noindent\textit{Source reference.} Original lemma 00TT: sources/stacks-examples-defos/pinned-original-source.tex, source lines 1195--1195. Cited by the retained statement or proof but not part of the collected spans; supplied only as a navigation pointer. The argument does not depend on its content.\par\medskip
\ifdefined\hypertarget\hypertarget{sourcepointer-2}{}\fi
\noindent\textit{Source reference.} Original lemma 00RY: sources/stacks-examples-defos/pinned-original-source.tex, source lines 1195--1195. Cited by the retained statement or proof but not part of the collected spans; supplied only as a navigation pointer. The argument does not depend on its content.\par\medskip
\ifdefined\hypertarget\hypertarget{sourcepointer-3}{}\fi
\noindent\textit{Source reference.} Original lemma 00RT: sources/stacks-examples-defos/pinned-original-source.tex, source lines 1195--1195. Cited by the retained statement or proof but not part of the collected spans; supplied only as a navigation pointer. The argument does not depend on its content.\par\medskip
\ifdefined\hypertarget\hypertarget{sourcepointer-4}{}\fi
\noindent\textit{Source reference.} Original lemma 0ALW: sources/stacks-examples-defos/pinned-original-source.tex, source lines 1195--1195. Cited by the retained statement or proof but not part of the collected spans; supplied only as a navigation pointer. The argument does not depend on its content.\par\medskip
\ifdefined\hypertarget\hypertarget{sourcepointer-5}{}\fi
\noindent\textit{Source reference.} Original lemma 0583: sources/stacks-examples-defos/pinned-original-source.tex, source lines 1195--1195. Cited by the retained statement or proof but not part of the collected spans; supplied only as a navigation pointer. The argument does not depend on its content.\par\medskip
\ifdefined\hypertarget\hypertarget{sourcepointer-6}{}\fi
\noindent\textit{Source reference.} Original lemma 00S7: sources/stacks-examples-defos/pinned-original-source.tex, source lines 1195--1195. Cited by the retained statement or proof but not part of the collected spans; supplied only as a navigation pointer. The argument does not depend on its content.\par\medskip
\ifdefined\hypertarget\hypertarget{sourcepointer-7}{}\fi
\noindent\textit{Source reference.} Original lemma 00J3: sources/stacks-examples-defos/pinned-original-source.tex, source lines 1195--1195. Cited by the retained statement or proof but not part of the collected spans; supplied only as a navigation pointer. The argument does not depend on its content.\par\medskip
\ifdefined\hypertarget\hypertarget{sourcepointer-8}{}\fi
\noindent\textit{Source reference.} Original lemma 07BU: sources/stacks-examples-defos/pinned-original-source.tex, source lines 1195--1195. Cited by the retained statement or proof but not part of the collected spans; supplied only as a navigation pointer. The argument does not depend on its content.\par\medskip
\ifdefined\hypertarget\hypertarget{sourcepointer-9}{}\fi
\noindent\textit{Source reference.} Original lemma 02CE: sources/stacks-examples-defos/pinned-original-source.tex, source lines 1195--1195. Cited by the retained statement or proof but not part of the collected spans; supplied only as a navigation pointer. The argument does not depend on its content.\par\medskip
\ifdefined\hypertarget\hypertarget{sourcepointer-10}{}\fi
\noindent\textit{Source reference.} Original lemma 00FV: sources/stacks-examples-defos/pinned-original-source.tex, source lines 1195--1195. Cited by the retained statement or proof but not part of the collected spans; supplied only as a navigation pointer. The argument does not depend on its content.\par\medskip
\ifdefined\hypertarget\hypertarget{sourcepointer-11}{}\fi
\noindent\textit{Source reference.} Original lemma 0DY3: sources/stacks-examples-defos/pinned-original-source.tex, source lines 1072--1079. Cited by the retained statement or proof but not part of the collected spans; supplied only as a navigation pointer. The argument does not depend on its content.\par\medskip
\ifdefined\hypertarget\hypertarget{sourcepointer-12}{}\fi
\noindent\textit{Source reference.} Original lemma 087R: sources/stacks-examples-defos/pinned-original-source.tex, source lines 1195--1195. Cited by the retained statement or proof but not part of the collected spans; supplied only as a navigation pointer. The argument does not depend on its content.\par\medskip
\begin{lemma}
\label{lemma-finite-type-rings-TI}
In Lemma \ref{lemma-rings-TI} if $P$ is a finite type $k$-algebra, then
\begin{enumerate}
\item $\text{Inf}(\Deformationcategory_P)$ is finite dimensional if and only if
$\dim(P) = 0$, and
\item $T\Deformationcategory_P$ is finite dimensional if
$\Spec(P) \to \Spec(k)$ is smooth except at a finite number of points.
\end{enumerate}
\end{lemma}

\begin{proof}
Proof of (1). We view $\text{Der}_k(P, P)$ as a $P$-module.
If it has finite dimension over $k$, then it has finite length
as a $P$-module, hence it is supported in finitely many
closed points of $\Spec(P)$
(Algebra, Lemma \ref{algebra-lemma-simple-pieces}).
Since $\text{Der}_k(P, P) = \Hom_P(\Omega_{P/k}, P)$
we see that
$\text{Der}_k(P, P)_\mathfrak p = \text{Der}_k(P_\mathfrak p, P_\mathfrak p)$
for any prime $\mathfrak p \subset P$
(this uses Algebra, Lemmas
\ref{algebra-lemma-differentials-localize},
\ref{algebra-lemma-differentials-finitely-presented}, and
\ref{algebra-lemma-hom-from-finitely-presented}).
Let $\mathfrak p$ be a minimal prime ideal of $P$
corresponding to an irreducible component of dimension $d > 0$.
Then $P_\mathfrak p$ is an Artinian local ring
essentially of finite type over $k$ with residue field
and $\Omega_{P_\mathfrak p/k}$ is nonzero for example by
Algebra, Lemma \ref{algebra-lemma-characterize-smooth-over-field}.
Any nonzero finite module over an Artinian local ring
has both a sub and a quotient module isomorphic to the residue field.
Thus we find that
$\text{Der}_k(P_\mathfrak p, P_\mathfrak p) =
\Hom_{P_\mathfrak p}(\Omega_{P_\mathfrak p/k}, P_\mathfrak p)$
is nonzero too. Combining all of the above we find that (1) is true.

\medskip\noindent
Proof of (2). For a prime $\mathfrak p$ of $P$ we will use that
$\NL_{P_\mathfrak p/k} = (\NL_{P/k})_\mathfrak p$
(Algebra, Lemma \ref{algebra-lemma-localize-NL})
and we will
use that
$\text{Ext}_P^1(\NL_{P/k}, P)_\mathfrak p =
\text{Ext}_{P_\mathfrak p}^1(\NL_{P_\mathfrak p/k}, P_\mathfrak p)$
(More on Algebra, Lemma
\ref{more-algebra-lemma-pseudo-coherence-and-base-change-ext}).
Given a prime $\mathfrak p \subset P$
then $k \to P$ is smooth at $\mathfrak p$ if and only if
$(\NL_{P/k})_\mathfrak p$ is quasi-isomorphic
to a finite projective module placed in degree $0$ (this follows
immediately from the definition of a smooth ring map but it also
follows from the stronger Algebra, Lemma \ref{algebra-lemma-smooth-at-point}).

\medskip\noindent
Assume that $P$ is smooth over $k$ at all but finitely many primes.
Then these ``bad'' primes are maximal ideals
$\mathfrak m_1, \ldots, \mathfrak m_n \subset P$ by
Algebra, Lemma \ref{algebra-lemma-finite-type-algebra-finite-nr-primes}
and the fact that the ``bad'' primes form a closed subset of $\Spec(P)$.
For $\mathfrak p \not \in \{\mathfrak m_1, \ldots, \mathfrak m_n\}$
we have $\text{Ext}^1_P(\NL_{P/k}, P)_\mathfrak p = 0$ by the results above.
Thus $\text{Ext}^1_P(\NL_{P/k}, P)$ is a finite $P$-module
whose support is contained in $\{\mathfrak m_1, \ldots, \mathfrak m_r\}$.
By Algebra, Proposition
\ref{algebra-proposition-minimal-primes-associated-primes}
for example, we find that the dimension over $k$ of
$\text{Ext}^1_P(\NL_{P/k}, P)$ is a finite integer combination
of $\dim_k \kappa(\mathfrak m_i)$ and hence finite by
the Hilbert Nullstellensatz
(Algebra, Theorem \ref{algebra-theorem-nullstellensatz}).
\end{proof}

\end{document}
